%% ===========================================================================
%%  Chapter 6 — Modular AI Core
%% ===========================================================================
\section{Modular AI Core}\label{ch:ai}

\texttt{assembly/sakum\_ai.s} is a from-scratch neural-ish engine in raw
x86-64.  It is the cognitive-ingestion half of the binary-hash system: each
\texttt{Knowledge/} chunk carries a \texttt{नाम hash = \#what <hex64>;}
line plus \texttt{नाम node = "..."} and \texttt{लेख(query("..."))} content,
and the core folds \emph{all three} into a 64-cell weight matrix
\texttt{W[]} so the forward pass reflects loaded knowledge rather than a
synthetic formula.

\subsection{Lifecycle}

\begin{enumerate}
  \item \texttt{walk\_dir} --- bounded 3-level directory walker.  At depth-1
    (\texttt{cat/}) and depth-2 (\texttt{cat/sub/}) it calls
    \texttt{try\_manifest} $\to$ \texttt{maybe\_load}; at depth-3 it
    string-matches the \texttt{manifest.sakum} entry and loads it.  Uses
    \texttt{opendir\$INODE64}/\texttt{readdir\$INODE64} (\texttt{d\_name} @
    offset 21, \texttt{d\_type} @ 20, \texttt{DT\_DIR=4}); \texttt{is\_dot}
    skips \texttt{.}/\texttt{..}, \texttt{is\_dir} skips non-directories.
  \item \texttt{maybe\_load(path)} --- opens via \texttt{\_open\$NOCANCEL},
    \textbf{fstat-verifies} the fd (Rosetta workaround, \S\ref{sec:rosetta}),
    \texttt{read}s the chunk into \texttt{rbuf}, calls \texttt{ingest\_chunk},
    then counts it under \texttt{budget}.
  \item \texttt{ingest\_chunk(buf, len)} --- scans for the \texttt{\#what }
    marker, parses the following hex digits and accumulates them into
    \texttt{W[]} modulo 9973, keyed by the running chunk index; then finds
    \texttt{node = "} and \texttt{query("} and folds those quoted strings
    into \texttt{W[]} as well.
  \item \texttt{load\_weights} + \texttt{forward\_pass} --- seeds the input
    vector \texttt{Vin}; \texttt{W[]} already carries the ingested
    knowledge.  Computes \texttt{Vout = W $\cdot$ Vin} (8$\times$8) and
    prints the first output.
  \item \texttt{self\_update} --- \texttt{fopen(ledger,"a")} + \texttt{fprintf}
    a tick line (see \S\ref{ch:datamodel}).
\end{enumerate}

\subsection{Data-flow}

\begin{verbatim}
Knowledge/ tree
      |  walk_dir (opendir$INODE64 / readdir$INODE64)
      v
manifest.sakum  --maybe_load-->  _open$NOCANCEL + _fstat$INODE64 (verify)
      |                                  |
      |  _read into rbuf                 v
      v                           ingest_chunk
   #what <hex>  ---->  W[idx] = (W[idx]*16 + digit) mod 9973
   node = "..."  --->  fold_str  ->  W[(base+i)%64]
   query("...")  --->  fold_str  ->  W[(base+32+i)%64]
      |
      v
   forward_pass:  Vout = W . Vin   (8x8)
      |
      v
   self_update -> ai_ledger.txt  ("ai tick: neurons=64 loaded=85 ...")
\end{verbatim}

\subsection{Ingestion: the binary-hash imprint}

\begin{lstlisting}[language=SakumAsm,caption={Hex fold into W (sakum\_ai.s:766).}]
.ic_hex:
    mov %r13, %rax
    mov $64, %r14
    xor %rdx, %rdx
    div %r14
    mov %rdx, %r14               # idx mod 64
    lea W(%rip), %r8
    mov (%r8,%r14,4), %eax       # cur weight
    imul $16, %eax
    movzx %dl, %edx
    add %edx, %eax
    xor %edx, %rdx
    mov $9973, %r9d
    div %r9d
    mov %edx, %eax              # mod 9973
    mov %eax, (%r8,%r14,4)
    inc %r13                    # advance weight index
\end{lstlisting}

\subsection{Content ingestion: substring + fold}

Two helpers make the content parse possible without a host language.

\textbf{\texttt{find\_sub}} searches a NUL-terminated needle in a buffer and
returns a pointer just past the first match (or 0):

\begin{lstlisting}[language=SakumAsm,caption={Substring search skeleton (sakum\_ai.s).}]
find_sub:                       # rdi=buf rsi=len rdx=needle -> rax=ptr|0
    push %rbx; push %rcx; push %r8; push %r9; push %r10; push %r11; push %r12
    ...                         # compute needle length, scan, compare
.fs_hit:
    lea (%r8,%rcx), %rax
    add %r11, %rax              # pointer just past needle
    pop %r12; ...; ret
.fs_none:
    xor %rax, %rax
    pop %r12; ...; ret
\end{lstlisting}

\textbf{\texttt{fold\_str}} folds the bytes of a quoted string into
\texttt{W[]} until a delimiter: \texttt{W[(base+i)\%64] = (W[...]*31 +
byte) mod 9973}.  The AI core calls it twice after the hash loop: once for
the node name (base = chunk index) and once for the query target (base =
chunk index + 32).

\begin{lstlisting}[language=SakumAsm,caption={Content scan in ingest\_chunk (sakum\_ai.s:789).}]
.ic_end:
    mov loaded(%rip), %r13          # base weight index = chunk number
    lea rbuf(%rip), %rdi
    mov %r15, %rsi
    lea s_node(%rip), %rdx          # needle "node = \""
    call find_sub
    test %rax, %rax
    jz .ic_q
    mov %rax, %rdi                  # start just after "node = \""
    mov $'"', %rsi                  # delimiter
    mov %r13, %rdx                  # base index
    call fold_str
.ic_q:
    lea rbuf(%rip), %rdi
    mov %r15, %rsi
    lea s_query(%rip), %rdx         # needle "query(\""
    call find_sub
    test %rax, %rax
    jz .ic_cdone
    ...                             # fold query at (chunk_idx+32)%64
    call fold_str
.ic_cdone:
\end{lstlisting}

\subsection{Forward pass}

\begin{lstlisting}[language=SakumAsm,caption={8x8 dot product (forward\_pass, sakum\_ai.s:860).}]
.fp_j:
    cmp $8, %r13d
    jge .fp_store
    mov %r12d, %ecx
    shl $3, %ecx
    add %r13d, %ecx
    lea W(%rip), %rbx
    mov (%rbx,%rcx,4), %edx         # W[i][j]
    lea Vin(%rip), %rbx
    mov (%rbx,%r13,4), %ecx         # Vin[j]
    imul %ecx, %edx
    add %edx, %eax
    inc %r13d
    jmp .fp_j
\end{lstlisting}

\subsection{Important: weights must not be overwritten}

An earlier bug had \texttt{load\_weights} running \emph{after}
\texttt{walk\_dir} and overwriting the entire \texttt{W[]} with a synthetic
formula, erasing all ingestion.  \texttt{load\_weights} now seeds only the
input vector \texttt{Vin}, leaving \texttt{W[]} intact.  With ingestion
active the forward-pass output is \texttt{out[0]=132256} (driven by the 85
chunks' hashes + names + queries); with the synthetic overwrite it was a
constant \texttt{37}.

\subsection{RAM-aware neuron scaling}

Neuron count grows by 8 per loaded chunk, capped at 64.  Total RAM is read
via \texttt{sysctl} (\texttt{CTL\_HW=6, HW\_MEMSIZE=24}) and reported in
MB in the ledger tick.

\subsection{Run proof}

\begin{verbatim}
$ gcc -arch x86_64 assembly/sakum_ai.s -o /tmp/ai && /tmp/ai
  inference out[0]=132256 (forward pass, 64 neurons)
  AI ready: 85 chunks loaded, 64 neurons active, 0 leaks.
\end{verbatim}
